%% ===================================================================
\section{Related work}
\label{sec:related}

The problem studied here lies at the intersection of three bodies of work:
propagation models and the digital twins that consume them for network
planning, learned and transferred propagation models, and distribution-free
uncertainty for decisions taken on predicted quantities. This section
reviews each in turn and closes with a feature-based comparison between the
representative studies and the present work.

\subsection{Propagation models and digital twins for network planning}
Empirical models such as Okumura--Hata~\cite{hata1980} and the 3GPP
TR~38.901 formulas~\cite{tr38901} express path loss through a small number
of parameters fitted to measurement campaigns and applied uniformly over an
environment class; TR~38.901 distinguishes a line-of-sight and a
non-line-of-sight law, which a digital twin can select per link from its
geometry. Measurement studies showed early that the path-loss exponent is
itself a random quantity that changes from cell to cell~\cite{erceg1999},
which is the variability a tile-level model tries to capture. Diffraction
over the buildings along a link is described by the knife-edge and
multiple-edge constructions of ITU-R P.526~\cite{itup526}, which serve in
this study as non-learned per-link baselines. Deterministic ray tracing
resolves the geometry explicitly~\cite{yun2015raytracing}; differentiable
ray tracers such as Sionna~RT make city-scale simulation and calibration
against measurements practical~\cite{hoydis2023sionna,hoydis2024learning}.
In this study, ray tracing produces the ground truth for many cities and is
not the run-time predictor.

Digital twins of wireless networks~\cite{khan2022dt,alkhateeb2023dt} and the
ITU-T digital twin network architecture~\cite{ituy3090} need propagation
models for regions where no measurements exist yet and consume them through
planning functions. Site selection with propagation predictions is a
classical optimisation problem~\cite{amaldi2003}; coverage maximisation is
submodular and admits greedy approximation guarantees~\cite{nemhauser1978}.
In operational networks, targeted measurements can be collected with the
minimisation-of-drive-tests (MDT) framework~\cite{johansson2012mdt,ts37320},
which motivates the few-shot setting of this study; choosing where to
measure is a submodular sensor-placement problem~\cite{krause2008}.

\subsection{Learned and transferred propagation models}
Convolutional~\cite{levie2021radiounet}, autoencoder~\cite{teganya2022},
adversarial~\cite{zhang2023rmegan} and diffusion~\cite{wang2025radiodiff}
models predict dense radio maps from environment maps, and radio-map
estimation from sparse measurements has a mature statistical
treatment~\cite{romero2022}. Channel knowledge maps~\cite{zeng2021ckm}
generalise the idea to other channel quantities. Regression of path loss
from engineered or image features has been studied with measurements at
\SI{2.6}{\giga\hertz}~\cite{thrane2020} and
\SI{28}{\giga\hertz}~\cite{gupta2022}, across several cities with
obstruction-height maps~\cite{dempsey2025}, and is surveyed
in~\cite{seretis2022}. These methods predict path loss or maps directly; the
tile-parameter approach studied here predicts a compact, physically
interpretable pair per tile that a planning tool can evaluate for any site
layout. The link model of this study belongs to the feature-regression
family of~\cite{thrane2020,dempsey2025}; its purpose is to serve as a
strong, simple and interpretable reference against which the tile
representation is measured.

MAPLE divides 3-D city maps into tiles and estimates mixed-path propagation
in a new deployment area from the similarity of its tiles to tiles of
measured areas~\cite{ozyurt2020maple}. The descriptor-based framework
of~\cite{ozyurt2026} describes each tile by 116 terrain, built-structure,
land-cover and climatic descriptors, calibrates log-distance parameters per
tile from measurements with the mixed-path composition of MAPLE, and
transfers them to unmeasured tiles by a Gaussian-kernel weighting of
descriptor similarity. It was evaluated with confidential measurements of an
industrial partner in two European cities, each serving once as the source
of the other, with tiles of \SI{200}{\metre} to \SI{1000}{\metre}, by the
normalised error of the received power, and against empirical models only.
Transfer between domains with different input distributions is covariate
shift~\cite{shimodaira2000}, and its difficulty depends on the discrepancy
between the domains~\cite{bendavid2010}; importance
weighting~\cite{shimodaira2000} and the alignment of feature
statistics~\cite{sun2016coral} are standard remedies, which are included
here as baselines.

\subsection{Uncertainty, censoring and positioning of this study}
Conformal prediction gives distribution-free prediction sets under
exchangeability~\cite{vovk2005}, and leave-one-out variants such as the
jackknife+ reuse every training point for calibration~\cite{barber2021};
this study applies the exchangeability argument at the level of cities, the
unit of shift. Measurements with a detection floor are censored, for which
the Tobit model is the classical treatment~\cite{tobin1958}; received-power
data beyond a receiver's sensitivity are censored in exactly this sense, and
a hurdle model that first predicts whether a link exists is the natural
counterpart for prediction.

Table~\ref{tab:related} summarises how the present study relates to the
representative works above. None of the learned propagation models has been
scored through the quantities a planner decides on, none states what
coverage can be promised from a transferred model, and the tile-transfer
framework has been evaluated on two cities with confidential data. The
present study addresses the same transfer problem on a larger public corpus
with a re-implementation of the kernel operator on the built-structure
descriptors, analyses that operator, compares it with learned regressors
under a common protocol, evaluates both the receiver-tile and the mixed-path
composition, measures the consequences for network planning, and states the
coverage that can be promised at a given risk.

\begin{table*}[!t]
\centering
\caption{Feature-based comparison between representative related studies
and the present work. Cross-city: trained on some cities and evaluated on
others. Tile parameters: a per-tile log-distance pair that a planning tool
can reuse. Link geometry: per-link obstruction or image features. Network
level: serving site, SINR, coverage or cell edges scored against the truth.
Planning: site selection or dimensioning. Guarantee: a statistical coverage
guarantee for the transferred model.}
\label{tab:related}
\footnotesize\setlength{\tabcolsep}{4pt}
\begin{tabular}{lcccccc>{\raggedright\arraybackslash}p{5.4cm}}
\toprule
Reference & Cross-city & Tile parameters & Link geometry & Network level &
Planning & Guarantee & Main limitation with respect to this study \\
\midrule
MAPLE~\cite{ozyurt2020maple} & \checkmark & \checkmark & -- & -- & -- & -- &
Tile-similarity estimation of mixed-path propagation; evaluated through the
propagation estimate rather than through planning quantities. \\
Ozyurt~\cite{ozyurt2026} & \checkmark & \checkmark & -- & -- & -- & -- &
Two cities with confidential measurements, each the source of the other;
received-power error against empirical models only. \\
Thrane et al.~\cite{thrane2020} & -- & -- & \checkmark & -- & -- & -- &
Measurements from one area at \SI{2.6}{\giga\hertz}; path loss only. \\
Gupta et al.~\cite{gupta2022} & -- & -- & \checkmark & -- & -- & -- &
One city at \SI{28}{\giga\hertz}; path loss only. \\
Dempsey et al.~\cite{dempsey2025} & \checkmark & -- & \checkmark & -- & -- & -- &
Path loss only; no reusable parameter representation and no planning
quantities. \\
RadioUNet~\cite{levie2021radiounet} & -- & -- & -- & -- & -- & -- &
Dense map prediction on simulated maps, split at random rather than held out
by city. \\
RadioDiff~\cite{wang2025radiodiff} & -- & -- & -- & -- & -- & -- &
Generative map prediction on simulated maps; no network-level evaluation. \\
Hoydis et al.~\cite{hoydis2024learning} & -- & -- & \checkmark & -- & -- & -- &
Calibration of a ray tracer to measurements of one area; the ray tracer
remains the predictor. \\
Amaldi et al.~\cite{amaldi2003} & -- & -- & -- & \checkmark & \checkmark & -- &
Optimisation with given propagation predictions; no learned or transferred
model. \\
\textbf{This work} & \checkmark & \checkmark & \checkmark & \checkmark &
\checkmark & \checkmark & Ray-traced rather than measured truth; eleven
cities limit the power of city-level tests. \\
\bottomrule
\end{tabular}
\end{table*}
